%% chapters/assignment/ch3_ue.tex
%% 第3章　利用者均衡配分（UE）

\section{活動経路の記述：アクティビティモデル}\label{sec:abms}

\sref{sec:ch1-network}で整備した記法と，\sref{sec:ch2-uncongested}で扱ったフロー非依存型配分を踏まえ，
本節では\textbf{利用者均衡配分}（User Equilibrium Assignment, UE）を体系的に扱う．
UE は，リンク旅行時間がリンク交通量に依存する混雑ネットワークにおいて，
利用者が自己の旅行時間を最小化しようとする合理的行動の均衡状態を記述するモデルであり，
交通量配分理論の根幹をなす\cite{Wardrop1952,Beckmann1956}．

本節の構成は次のとおりである．
\ssref{subsec:wardrop} で均衡条件を与える Wardrop の第1原則を導入し，
\ssref{subsec:beckmann} でこの条件を等価な凸最適化問題（Beckmann 変換）に変換する．
\ssref{subsec:frank-wolfe} では，この最適化問題の標準的な解法である
\textbf{Frank-Wolfe 法}を他手法との比較も交えて説明する．
\ssref{subsec:braess} では，UE の本質的な性質を浮き彫りにする事例として
\textbf{Braess のパラドックス}を紹介する．

% ---------------------------------------------------------------
\subsection{Wardrop の第1原則}
\label{subsec:wardrop}
% ---------------------------------------------------------------

\subsubsection*{均衡条件の直感的背景}

混雑を考慮したネットワークでは，リンク旅行時間 $t_a(x_a)$ は
リンク交通量 $x_a$ の増加関数である（\sref{sec:ch1-network} \ssref{subsec:bpr} 参照）．
このとき，多数の利用者が同時に自己の旅行時間を最小化しようとすれば，
最初は最短経路に交通量が集中して旅行時間が上昇し，
やがて各利用者が経路をそれ以上変更する動機を失う状態